\documentclass[12pt, a4paper]{article}
\usepackage[utf8]{inputenc}
\usepackage[T2A]{fontenc} 
\usepackage[russian]{babel}
\usepackage{amsmath, amssymb}
\usepackage{graphicx}
\usepackage{tikz}
\usetikzlibrary{patterns, arrows.meta}
\usepackage{float}
\usepackage{array}

\usepackage[left=3cm, right=1.5cm, top=2cm, bottom=2cm]{geometry}

\setlength{\parindent}{0pt}
\setlength{\parskip}{1em}
\setlength{\headheight}{15pt}
\addtolength{\topmargin}{-3pt}
\usepackage{fancyhdr}
\pagestyle{fancy}
\fancyhf{} 
\fancyhead[C]{\thepage} 
\renewcommand{\headrulewidth}{0pt} 

\usepackage{titlesec}
\titleformat{\section}
  {\centering\Large\bfseries} 
  {\thesection}               
  {1em}                       
  {}

\titleformat{\subsection}
  {\centering\large\bfseries}
  {\thesubsection}            
  {1em}                       
  {}

\begin{document}

\thispagestyle{empty}

\begin{center}
    \textbf{Бюджетное учреждение высшего образования\\
    Ханты-Мансийского автономного округа -- Югры\\
    «СУРГУТСКИЙ ГОСУДАРСТВЕННЫЙ УНИВЕРСИТЕТ»}\\[0.8cm]
    
    \textbf{Политехнический институт}\\[0.1cm]
    \rule{\textwidth}{1.5pt}\\[0.6cm]
    
    Кафедра прикладной математики\\
    \text{Гркикян М.Э.}\\[1.2cm]
    
    \textbf{Численное решение уравнения пьезопроводности \\ 
    с помощью физически информированных нейронных сетей}\\[1.2cm]
    
    БАКАЛАВРСКАЯ РАБОТА\\
    по направлению 01.03.02 «Прикладная математика и информатика»\\
    направленность (профиль): «Прикладная математика и информатика»
\end{center}

\vspace{1.5cm}

\begin{flushright}
    \begin{tabular}{l}
        Научный руководитель:\\
        Зав. кафедрой ПМ,\\
        доцент, к.ф.-м.н. \text{Гореликов А.В.}\\[0.2cm]
        \rule{5cm}{0.5pt}\\
        \footnotesize (подпись)\\[1.5cm]
        
        Допущено к защите:\\
        \rule{2.5cm}{0.5pt} 2026 г.\\
        Зав. кафедрой ПМ,\\
        доцент, к.ф.-м.н. Гореликов А.В.\\[0.2cm]
        \rule{5cm}{0.5pt}\\
        \footnotesize (подпись)\\[1.5cm]
        
        Студент гр. № 601-31\\
        \text{Гркикян М.Э.}\\[0.2cm]
        \rule{5cm}{0.5pt}\\
        \footnotesize (подпись)
    \end{tabular}
\end{flushright}

\vfill % Прижимаем город и год к самому низу страницы

\begin{center}
    Сургут 2026 г.
\end{center}

\newpage
\tableofcontents
\newpage

\section*{СПИСОК ОБОЗНАЧЕНИЙ}
\addcontentsline{toc}{section}{СПИСОК ОБОЗНАЧЕНИЙ}

\begin{table}[H]
    \centering
    \renewcommand{\arraystretch}{1.5} % Чуть увеличиваем высоту строк для красоты
    \begin{tabular}{|>{\raggedright\arraybackslash}p{7cm}|c|>{\raggedright\arraybackslash}p{4cm}|}
        \hline
        \textbf{Величина} & \textbf{Обозначение} & \textbf{Единицы измерения} \\ \hline
        Давление & \(p\) & Па \\ \hline
        Начальное давление & \(p^0\) & Па \\ \hline
        Время & \(t\) & с \\ \hline
        Координата & \(x\) & м \\ \hline
        Длина расчетного участка & \(L\) & м \\ \hline
        Длина элементарного объема & \(\Delta x\) & м \\ \hline
        Скорость фильтрации & \(v\) & м/с \\ \hline
        Проницаемость породы & \(k\) & м\(^2\) \\ \hline
        Динамическая вязкость жидкости & \(\mu\) & Па\(\cdot\)с \\ \hline
        Плотность жидкости & \(\rho\) & кг/м\(^3\) \\ \hline
        Начальная плотность жидкости & \(\rho^0\) & кг/м\(^3\) \\ \hline
        Пористость породы & \(\phi\) & безразмерная \\ \hline
        Начальная пористость породы & \(\phi^0\) & безразмерная \\ \hline
        Убывание массы в единице объема & \(\tilde{q}\) & кг/(м\(^3\cdot\)с) \\ \hline
        Сжимаемость жидкости & \(c_f\) & Па\(^{-1}\) \\ \hline
        Сжимаемость породы & \(c_R\) & Па\(^{-1}\) \\ \hline
        Коэффициент пьезопроводности & \(\kappa\) & м\(^2\)/с \\ \hline
    \end{tabular}
\end{table}

\newpage

\section*{ВВЕДЕНИЕ}
\addcontentsline{toc}{section}{ВВЕДЕНИЕ}

\textbf{Актуальность работы.} Уравнение пьезопроводности описывает неустановившуюся фильтрацию упругой жидкости в пористой среде и лежит в основе многих задач разработки нефтегазовых месторождений, гидрогеологии и подземной гидромеханики.

В последние годы активно развивается альтернативный подход — физически информированные нейронные сети (Physics-Informed Neural Networks, PINN).

\textbf{Цель работы.} Реализовать и исследовать физически информированную нейронную сеть для численного решения уравнения пьезопроводности в одномерной постановке.

\textbf{Задачи работы:}
\begin{enumerate}
    \item Изучить математическую постановку уравнения пьезопроводности, типовые начальные и граничные условия.
    \item Провести обзор существующих примеров применения PINN к уравнениям параболического типа.
    \item Реализовать полносвязную нейронную сеть и настроить автоматическое дифференцирование для вычисления производных, входящих в уравнение.
    \item Сформировать функцию потерь с учетом невязки уравнения, начального и граничных условий.
    \item Обучить сеть и провести серию вычислительных экспериментов.
    \item Сравнить полученные решения с аналитическими и конечно-разностными результатами.
\end{enumerate}

\textbf{Объект исследования.} Процесс неустановившейся одномерной фильтрации упругой жидкости в пористой среде.

\textbf{Предмет исследования.} Решение уравнения пьезопроводности при помощи физически информированных нейронных сетей (PINN).

\textbf{Методы исследования.}

\textbf{Практическая значимость.} Разработанный код и методика применения PINN могут быть использованы в качестве безсеточного инструмента для решения простых задач фильтрации.

\textbf{Структура работы.} Работа состоит из введения, четырех глав, заключения и списка литературы.

\newpage
\section{Краткий обзор литературы по теме ВКР}

\newpage

\section{Математическая постановка задачи}

Рассматривается одномерная неустановившаяся фильтрация упругой жидкости в пористой среде.

В основе математической модели лежат закон Дарси и уравнение сохранения массы.

Закон Дарси в одномерном случае имеет вид:
\begin{equation}
    v = -\frac{k}{\mu} \frac{\partial p}{\partial x},
\end{equation}
где \(v\) — скорость фильтрации, \(k\) — проницаемость породы, \(\mu\) — вязкость жидкости, \(p\) — давление.

Для вывода уравнения сохранения массы рассмотрим элементарный объем пористой среды (рис. \ref{fig:scheme}).

\begin{figure}[H]
    \centering
    \begin{tikzpicture}[>=Stealth, thick, scale=0.85, transform shape]
        % Основной цилиндр
        \draw (0,0) rectangle (8, 3);
        
        % Ось x
        \draw[->] (1, 1.5) -- (4, 1.5) node[midway, above] {\(x\)};
        
        % Длина L
        \draw[<->] (0, 3.8) -- (8, 3.8) node[midway, above] {\(L\)};
        \draw[dashed] (0,3) -- (0,3.8);
        \draw[dashed] (8,3) -- (8,3.8);
        
        % Элементарный объем (штриховка)
        \draw[pattern=crosshatch, pattern color=black] (4.5,0) rectangle (5.5,3);
        
        % Дельта x
        \draw[<->] (4.5, -0.6) -- (5.5, -0.6) node[midway, below] {\(\Delta x\)};
        \draw[dashed] (4.5,0) -- (4.5,-0.6);
        \draw[dashed] (5.5,0) -- (5.5,-0.6);
        
        % Подпись к элементу (выносим наружу)
        \draw (5.5, 0.5) -- (6.8, -0.6) node[right, align=left] {Элементарный\\объем};
        
        % Стрелки притока (слева)
        \draw[->] (-1.5, 0.5) -- (0, 0.5);
        \draw[->] (-1.5, 1.5) -- (0, 1.5);
        \draw[->] (-1.5, 2.5) -- (0, 2.5);
        \node[left] at (-1.5, 1.5) {Приток};
        
        % Стрелки оттока (справа)
        \draw[->] (8, 0.5) -- (9.5, 0.5);
        \draw[->] (8, 1.5) -- (9.5, 1.5);
        \draw[->] (8, 2.5) -- (9.5, 2.5);
        \node[right] at (9.5, 1.5) {Отток};
    \end{tikzpicture}
    \caption{Схема линейной фильтрации в цилиндрическом образце длины \(\Delta x\)}
    \label{fig:scheme}
\end{figure}

Уравнение сохранения массы:
\begin{equation}
    \frac{\partial (\rho \phi)}{\partial t} + \frac{\partial (\rho v)}{\partial x} = - \tilde{q},
\end{equation}
где \(\tilde{q}\) — убывание массы в единице объема в единицу времени (положительно при отборе и отрицательно при нагнетании).

Подставим закон Дарси в уравнение сохранения массы:
\begin{equation}
    \frac{\partial (\rho \phi)}{\partial t} = \frac{\partial}{\partial x} \left( \rho \frac{k}{\mu} \frac{\partial p}{\partial x} \right) - \tilde{q}.
\end{equation}

Для вывода уравнения пьезопроводности используем линеаризованное уравнение состояния упругой жидкости и соотношение, описывающее изменение пористости в зависимости от давления:
\begin{equation}
    \rho = \rho^0 [1 + c_f (p - p^0)], \quad \phi = \phi^0 [1 + c_R (p - p^0)],
\end{equation}
где \(c_f\) и \(c_R\) — сжимаемости жидкости и породы, \(\rho^0\) и \(\phi^0\) — значения при начальном давлении \(p^0\).

Тогда производные по времени равны:
\begin{equation}
    \frac{\partial \rho}{\partial t} = \rho^0 c_f \frac{\partial p}{\partial t}, \quad \frac{\partial \phi}{\partial t} = \phi^0 c_R \frac{\partial p}{\partial t}.
\end{equation}

Раскроем производную в левой части уравнения:
\begin{equation}
    \frac{\partial (\rho \phi)}{\partial t} = \phi \frac{\partial \rho}{\partial t} + \rho \frac{\partial \phi}{\partial t}.
\end{equation}
Подставим уравнения состояния в раскрытую производную:
\begin{equation}
    \frac{\partial (\rho \phi)}{\partial t} = \phi^0 [1 + c_R (p - p^0)] \cdot \rho^0 c_f \frac{\partial p}{\partial t} + \rho^0 [1 + c_f (p - p^0)] \cdot \phi^0 c_R \frac{\partial p}{\partial t}.
\end{equation}

Вынесем общий множитель \(\rho^0 \phi^0 \frac{\partial p}{\partial t}\) и раскроем скобки:
\begin{equation}
    \frac{\partial (\rho \phi)}{\partial t} = \rho^0 \phi^0 [c_f + c_f c_R (p - p^0) + c_R + c_f c_R (p - p^0)] \frac{\partial p}{\partial t}.
\end{equation}

Приведем подобные слагаемые:
\begin{equation}
    \frac{\partial (\rho \phi)}{\partial t} = \rho^0 \phi^0 [c_f + c_R + 2 c_f c_R (p - p^0)] \frac{\partial p}{\partial t}.
\end{equation}

Так как \(c_f\) и \(c_R\) имеют порядок \(10^{-5}\) -- \(10^{-6}\), их произведение \(c_f \cdot c_R\) имеет порядок \(10^{-10}\) -- \(10^{-12}\). Пренебрегая членом \(2 c_f c_R (p - p^0)\), приходим к выражению:
\begin{equation}
    \frac{\partial (\rho \phi)}{\partial t} \approx \rho^0 \phi^0 (c_f + c_R) \frac{\partial p}{\partial t}.
\end{equation}

Теперь упростим правую часть. Так как сжимаемость мала, плотность в множителе перед производной по координате можно приближенно заменить на \(\rho^0\). Проницаемость \(k\) и вязкость \(\mu\) считаем постоянными:
\begin{equation}
    \frac{\partial}{\partial x} \left( \rho \frac{k}{\mu} \frac{\partial p}{\partial x} \right) \approx \rho^0 \frac{k}{\mu} \frac{\partial^2 p}{\partial x^2}.
\end{equation}

Приравняем упрощенные левую и правую части:
\begin{equation}
    \rho^0 \phi^0 (c_f + c_R) \frac{\partial p}{\partial t} = \rho^0 \frac{k}{\mu} \frac{\partial^2 p}{\partial x^2} - \tilde{q}.
\end{equation}

Разделим обе части на \(\rho^0 \phi^0 (c_f + c_R)\) и введем обозначение для коэффициента пьезопроводности:
\begin{equation}
    \frac{\partial p}{\partial t} = \kappa \frac{\partial^2 p}{\partial x^2} - \frac{\tilde{q}}{\rho^0 \phi^0 (c_f + c_R)},
\end{equation}
где
\begin{equation}
    \kappa = \frac{k}{\mu \phi^0 (c_f + c_R)}.
\end{equation}

Будем рассматривать случай, где \(\tilde{q} = 0\). Получаем классический вид уравнения пьезопроводности:

\begin{equation}
    \frac{\partial p}{\partial t} = \kappa \frac{\partial^2 p}{\partial x^2}.
\end{equation}

Для получения единственного решения уравнения пьезопроводности необходимо задать начальное и граничные условия.

Начальное условие задает давление в начальный момент времени \(t=0\). Будем считать, что в начальный момент времени давление в пласте нулевое:
\begin{equation}
    p(x, 0) = 0.
\end{equation}

Граничные условия задают поведение давления на концах рассматриваемого участка, например \(x = 0\) и \(x = 1\). В теории дифференциальных уравнений выделяют три основных типа граничных условий: условие Дирихле (задается само значение искомой функции), условие Неймана (задается производная искомой функции) и условие Робена (задается линейная комбинация самого значения функции и ее производной). 

В данной работе в качестве граничных используем условия первого рода (Дирихле), задающие фиксированные значения давления на границах:
\begin{equation}
    p(0, t) = 1, \quad p(1, t) = 0.
\end{equation}

Таким образом, итоговая математическая постановка одномерной задачи пьезопроводности принимает вид:
\begin{equation}
    \begin{cases}
        \frac{\partial p}{\partial t} = \kappa \frac{\partial^2 p}{\partial x^2} & 0 < x < 1, \quad t > 0, \\
        p(x, 0) = 0 & 0 \le x \le 1, \\
        p(0, t) = 1, \quad p(1, t) = 0 & t > 0.
    \end{cases}
\end{equation}

\newpage
\section{Численные методы и программная реализация}

\newpage
\section{Результаты расчётов и их анализ}

% Решение данной задачи описывает процесс перераспределения давления в пласте. В начальный момент времени давление внутри пласта равно нулю. За счет разности давлений на границах (слева \(p=1\), справа \(p=0\)) начинается фильтрация жидкости. С течением времени профиль давления эволюционирует и выходит на стационарный линейный режим. Эволюция профиля давления схематично представлена на рисунке \ref{fig:pressure_evolution}.

\begin{figure}[H]
    \centering
    \begin{tikzpicture}[x=7cm, y=5cm, >=Stealth, thick]
\draw[step=0.2, gray!20, very thin] (0,0) grid (1,1);
\foreach \x in {0.2, 0.4, 0.6, 0.8, 1} {
  \draw (\x, 2pt) -- (\x, -2pt) node[below] {\x};
}
\foreach \y in {0.2, 0.4, 0.6, 0.8, 1} {
  \draw (2pt, \y) -- (-2pt, \y) node[left] {\y};
}
\node[below left] at (0,0) {0};
\draw[->, thick] (0,0) -- (1.15,0) node[below right] {\(x\)};
\draw[->, thick] (0,0) -- (0,1.15) node[left] {\(p\)};
\draw[blue, opacity=0.4] plot coordinates {(0.000,1.000)(0.020,0.919)(0.041,0.839)(0.061,0.760)(0.082,0.684)(0.102,0.611)(0.122,0.541)(0.143,0.476)(0.163,0.415)(0.184,0.359)(0.204,0.308)(0.224,0.262)(0.245,0.221)(0.265,0.185)(0.286,0.154)(0.306,0.126)(0.327,0.103)(0.347,0.083)(0.367,0.066)(0.388,0.053)(0.408,0.041)(0.429,0.032)(0.449,0.025)(0.469,0.019)(0.490,0.014)(0.510,0.011)(0.531,0.008)(0.551,0.006)(0.571,0.004)(0.592,0.003)(0.612,0.002)(0.633,0.001)(0.653,0.001)(0.673,0.001)(0.694,0.000)(0.714,0.000)(0.735,0.000)(0.755,0.000)(0.776,0.000)(0.796,0.000)(0.816,0.000)(0.837,0.000)(0.857,0.000)(0.878,0.000)(0.898,0.000)(0.918,0.000)(0.939,0.000)(0.959,0.000)(0.980,0.000)(1.000,0.000)};
\draw[blue, opacity=0.4] plot coordinates {(0.000,1.000)(0.020,0.943)(0.041,0.885)(0.061,0.829)(0.082,0.773)(0.102,0.719)(0.122,0.665)(0.143,0.614)(0.163,0.564)(0.184,0.517)(0.204,0.471)(0.224,0.428)(0.245,0.387)(0.265,0.349)(0.286,0.313)(0.306,0.279)(0.327,0.249)(0.347,0.220)(0.367,0.194)(0.388,0.171)(0.408,0.149)(0.429,0.130)(0.449,0.113)(0.469,0.097)(0.490,0.083)(0.510,0.071)(0.531,0.061)(0.551,0.051)(0.571,0.043)(0.592,0.036)(0.612,0.030)(0.633,0.025)(0.653,0.021)(0.673,0.017)(0.694,0.014)(0.714,0.011)(0.735,0.009)(0.755,0.008)(0.776,0.006)(0.796,0.005)(0.816,0.004)(0.837,0.003)(0.857,0.002)(0.878,0.002)(0.898,0.001)(0.918,0.001)(0.939,0.001)(0.959,0.000)(0.980,0.000)(1.000,0.000)};
\draw[blue, opacity=0.4] plot coordinates {(0.000,1.000)(0.020,0.953)(0.041,0.906)(0.061,0.860)(0.082,0.814)(0.102,0.768)(0.122,0.724)(0.143,0.680)(0.163,0.638)(0.184,0.596)(0.204,0.556)(0.224,0.517)(0.245,0.480)(0.265,0.444)(0.286,0.410)(0.306,0.377)(0.327,0.346)(0.347,0.317)(0.367,0.289)(0.388,0.263)(0.408,0.239)(0.429,0.216)(0.449,0.195)(0.469,0.176)(0.490,0.157)(0.510,0.141)(0.531,0.126)(0.551,0.112)(0.571,0.099)(0.592,0.088)(0.612,0.077)(0.633,0.068)(0.653,0.059)(0.673,0.052)(0.694,0.045)(0.714,0.039)(0.735,0.034)(0.755,0.029)(0.776,0.025)(0.796,0.021)(0.816,0.018)(0.837,0.015)(0.857,0.012)(0.878,0.010)(0.898,0.008)(0.918,0.006)(0.939,0.005)(0.959,0.003)(0.980,0.001)(1.000,0.000)};
\draw[blue, opacity=0.4] plot coordinates {(0.000,1.000)(0.020,0.959)(0.041,0.919)(0.061,0.878)(0.082,0.838)(0.102,0.799)(0.122,0.760)(0.143,0.721)(0.163,0.683)(0.184,0.646)(0.204,0.610)(0.224,0.575)(0.245,0.541)(0.265,0.507)(0.286,0.475)(0.306,0.444)(0.327,0.414)(0.347,0.386)(0.367,0.359)(0.388,0.332)(0.408,0.308)(0.429,0.284)(0.449,0.262)(0.469,0.241)(0.490,0.221)(0.510,0.202)(0.531,0.185)(0.551,0.168)(0.571,0.153)(0.592,0.139)(0.612,0.125)(0.633,0.113)(0.653,0.102)(0.673,0.091)(0.694,0.082)(0.714,0.073)(0.735,0.065)(0.755,0.057)(0.776,0.050)(0.796,0.044)(0.816,0.038)(0.837,0.033)(0.857,0.028)(0.878,0.023)(0.898,0.019)(0.918,0.015)(0.939,0.011)(0.959,0.007)(0.980,0.004)(1.000,0.000)};
\draw[blue, opacity=0.4] plot coordinates {(0.000,1.000)(0.020,0.964)(0.041,0.927)(0.061,0.891)(0.082,0.855)(0.102,0.820)(0.122,0.784)(0.143,0.749)(0.163,0.715)(0.184,0.681)(0.204,0.648)(0.224,0.616)(0.245,0.584)(0.265,0.553)(0.286,0.523)(0.306,0.494)(0.327,0.465)(0.347,0.438)(0.367,0.411)(0.388,0.386)(0.408,0.361)(0.429,0.338)(0.449,0.315)(0.469,0.293)(0.490,0.273)(0.510,0.253)(0.531,0.234)(0.551,0.217)(0.571,0.200)(0.592,0.184)(0.612,0.169)(0.633,0.155)(0.653,0.142)(0.673,0.129)(0.694,0.117)(0.714,0.106)(0.735,0.096)(0.755,0.086)(0.776,0.077)(0.796,0.068)(0.816,0.060)(0.837,0.052)(0.857,0.045)(0.878,0.038)(0.898,0.031)(0.918,0.024)(0.939,0.018)(0.959,0.012)(0.980,0.006)(1.000,0.000)};
\draw[blue, opacity=0.4] plot coordinates {(0.000,1.000)(0.020,0.967)(0.041,0.934)(0.061,0.901)(0.082,0.868)(0.102,0.835)(0.122,0.803)(0.143,0.771)(0.163,0.739)(0.184,0.708)(0.204,0.677)(0.224,0.647)(0.245,0.617)(0.265,0.588)(0.286,0.559)(0.306,0.532)(0.327,0.505)(0.347,0.478)(0.367,0.453)(0.388,0.428)(0.408,0.404)(0.429,0.380)(0.449,0.358)(0.469,0.336)(0.490,0.315)(0.510,0.295)(0.531,0.276)(0.551,0.258)(0.571,0.240)(0.592,0.223)(0.612,0.207)(0.633,0.191)(0.653,0.177)(0.673,0.162)(0.694,0.149)(0.714,0.136)(0.735,0.124)(0.755,0.112)(0.776,0.101)(0.796,0.090)(0.816,0.080)(0.837,0.070)(0.857,0.061)(0.878,0.051)(0.898,0.042)(0.918,0.034)(0.939,0.025)(0.959,0.017)(0.980,0.008)(1.000,0.000)};
\draw[blue, opacity=0.4] plot coordinates {(0.000,1.000)(0.020,0.969)(0.041,0.938)(0.061,0.908)(0.082,0.877)(0.102,0.847)(0.122,0.817)(0.143,0.787)(0.163,0.757)(0.184,0.728)(0.204,0.699)(0.224,0.671)(0.245,0.643)(0.265,0.615)(0.286,0.588)(0.306,0.562)(0.327,0.536)(0.347,0.510)(0.367,0.486)(0.388,0.461)(0.408,0.438)(0.429,0.415)(0.449,0.393)(0.469,0.371)(0.490,0.350)(0.510,0.330)(0.531,0.311)(0.551,0.292)(0.571,0.273)(0.592,0.256)(0.612,0.239)(0.633,0.222)(0.653,0.206)(0.673,0.191)(0.694,0.176)(0.714,0.162)(0.735,0.148)(0.755,0.135)(0.776,0.122)(0.796,0.110)(0.816,0.098)(0.837,0.086)(0.857,0.074)(0.878,0.063)(0.898,0.052)(0.918,0.042)(0.939,0.031)(0.959,0.021)(0.980,0.010)(1.000,0.000)};
\draw[blue, opacity=0.4] plot coordinates {(0.000,1.000)(0.020,0.971)(0.041,0.942)(0.061,0.913)(0.082,0.885)(0.102,0.856)(0.122,0.828)(0.143,0.800)(0.163,0.772)(0.184,0.744)(0.204,0.717)(0.224,0.690)(0.245,0.663)(0.265,0.637)(0.286,0.611)(0.306,0.586)(0.327,0.561)(0.347,0.536)(0.367,0.512)(0.388,0.489)(0.408,0.466)(0.429,0.443)(0.449,0.421)(0.469,0.400)(0.490,0.379)(0.510,0.359)(0.531,0.339)(0.551,0.320)(0.571,0.301)(0.592,0.283)(0.612,0.265)(0.633,0.248)(0.653,0.231)(0.673,0.215)(0.694,0.199)(0.714,0.184)(0.735,0.169)(0.755,0.154)(0.776,0.140)(0.796,0.126)(0.816,0.113)(0.837,0.099)(0.857,0.086)(0.878,0.074)(0.898,0.061)(0.918,0.049)(0.939,0.036)(0.959,0.024)(0.980,0.012)(1.000,0.000)};
\draw[blue, opacity=0.4] plot coordinates {(0.000,1.000)(0.020,0.973)(0.041,0.945)(0.061,0.918)(0.082,0.891)(0.102,0.864)(0.122,0.837)(0.143,0.810)(0.163,0.784)(0.184,0.757)(0.204,0.731)(0.224,0.705)(0.245,0.680)(0.265,0.655)(0.286,0.630)(0.306,0.605)(0.327,0.581)(0.347,0.557)(0.367,0.534)(0.388,0.511)(0.408,0.488)(0.429,0.466)(0.449,0.445)(0.469,0.423)(0.490,0.403)(0.510,0.382)(0.531,0.362)(0.551,0.343)(0.571,0.324)(0.592,0.305)(0.612,0.287)(0.633,0.269)(0.653,0.252)(0.673,0.235)(0.694,0.218)(0.714,0.202)(0.735,0.186)(0.755,0.170)(0.776,0.155)(0.796,0.140)(0.816,0.125)(0.837,0.111)(0.857,0.096)(0.878,0.082)(0.898,0.068)(0.918,0.054)(0.939,0.041)(0.959,0.027)(0.980,0.014)(1.000,0.000)};
\draw[blue, opacity=0.4] plot coordinates {(0.000,1.000)(0.020,0.974)(0.041,0.948)(0.061,0.922)(0.082,0.896)(0.102,0.870)(0.122,0.844)(0.143,0.819)(0.163,0.793)(0.184,0.768)(0.204,0.743)(0.224,0.718)(0.245,0.693)(0.265,0.669)(0.286,0.645)(0.306,0.621)(0.327,0.598)(0.347,0.575)(0.367,0.552)(0.388,0.529)(0.408,0.507)(0.429,0.485)(0.449,0.464)(0.469,0.443)(0.490,0.422)(0.510,0.401)(0.531,0.381)(0.551,0.362)(0.571,0.342)(0.592,0.323)(0.612,0.305)(0.633,0.287)(0.653,0.269)(0.673,0.251)(0.694,0.234)(0.714,0.217)(0.735,0.200)(0.755,0.184)(0.776,0.167)(0.796,0.151)(0.816,0.136)(0.837,0.120)(0.857,0.105)(0.878,0.089)(0.898,0.074)(0.918,0.059)(0.939,0.044)(0.959,0.030)(0.980,0.015)(1.000,0.000)};
\draw[blue, opacity=0.4] plot coordinates {(0.000,1.000)(0.020,0.975)(0.041,0.950)(0.061,0.925)(0.082,0.900)(0.102,0.875)(0.122,0.850)(0.143,0.826)(0.163,0.801)(0.184,0.777)(0.204,0.752)(0.224,0.728)(0.245,0.705)(0.265,0.681)(0.286,0.658)(0.306,0.634)(0.327,0.611)(0.347,0.589)(0.367,0.566)(0.388,0.544)(0.408,0.522)(0.429,0.501)(0.449,0.479)(0.469,0.458)(0.490,0.438)(0.510,0.417)(0.531,0.397)(0.551,0.377)(0.571,0.358)(0.592,0.339)(0.612,0.320)(0.633,0.301)(0.653,0.283)(0.673,0.265)(0.694,0.247)(0.714,0.229)(0.735,0.212)(0.755,0.194)(0.776,0.178)(0.796,0.161)(0.816,0.144)(0.837,0.128)(0.857,0.111)(0.878,0.095)(0.898,0.079)(0.918,0.063)(0.939,0.047)(0.959,0.032)(0.980,0.016)(1.000,0.000)};
\draw[blue, opacity=0.4] plot coordinates {(0.000,1.000)(0.020,0.976)(0.041,0.952)(0.061,0.927)(0.082,0.903)(0.102,0.879)(0.122,0.855)(0.143,0.831)(0.163,0.807)(0.184,0.784)(0.204,0.760)(0.224,0.737)(0.245,0.714)(0.265,0.691)(0.286,0.668)(0.306,0.645)(0.327,0.623)(0.347,0.600)(0.367,0.578)(0.388,0.556)(0.408,0.535)(0.429,0.513)(0.449,0.492)(0.469,0.471)(0.490,0.451)(0.510,0.430)(0.531,0.410)(0.551,0.390)(0.571,0.371)(0.592,0.351)(0.612,0.332)(0.633,0.313)(0.653,0.294)(0.673,0.276)(0.694,0.257)(0.714,0.239)(0.735,0.221)(0.755,0.203)(0.776,0.186)(0.796,0.168)(0.816,0.151)(0.837,0.134)(0.857,0.117)(0.878,0.100)(0.898,0.083)(0.918,0.067)(0.939,0.050)(0.959,0.033)(0.980,0.017)(1.000,0.000)};
\draw[blue, opacity=0.4] plot coordinates {(0.000,1.000)(0.020,0.976)(0.041,0.953)(0.061,0.929)(0.082,0.906)(0.102,0.883)(0.122,0.859)(0.143,0.836)(0.163,0.813)(0.184,0.790)(0.204,0.767)(0.224,0.744)(0.245,0.721)(0.265,0.698)(0.286,0.676)(0.306,0.654)(0.327,0.632)(0.347,0.610)(0.367,0.588)(0.388,0.566)(0.408,0.545)(0.429,0.524)(0.449,0.503)(0.469,0.482)(0.490,0.461)(0.510,0.441)(0.531,0.421)(0.551,0.401)(0.571,0.381)(0.592,0.361)(0.612,0.342)(0.633,0.323)(0.653,0.304)(0.673,0.285)(0.694,0.266)(0.714,0.248)(0.735,0.229)(0.755,0.211)(0.776,0.193)(0.796,0.175)(0.816,0.157)(0.837,0.139)(0.857,0.122)(0.878,0.104)(0.898,0.087)(0.918,0.069)(0.939,0.052)(0.959,0.035)(0.980,0.017)(1.000,0.000)};
\draw[blue, opacity=0.4] plot coordinates {(0.000,1.000)(0.020,0.977)(0.041,0.954)(0.061,0.931)(0.082,0.908)(0.102,0.885)(0.122,0.862)(0.143,0.840)(0.163,0.817)(0.184,0.794)(0.204,0.772)(0.224,0.749)(0.245,0.727)(0.265,0.705)(0.286,0.683)(0.306,0.661)(0.327,0.639)(0.347,0.617)(0.367,0.596)(0.388,0.575)(0.408,0.553)(0.429,0.532)(0.449,0.511)(0.469,0.491)(0.490,0.470)(0.510,0.450)(0.531,0.429)(0.551,0.409)(0.571,0.389)(0.592,0.370)(0.612,0.350)(0.633,0.331)(0.653,0.311)(0.673,0.292)(0.694,0.273)(0.714,0.254)(0.735,0.236)(0.755,0.217)(0.776,0.198)(0.796,0.180)(0.816,0.162)(0.837,0.144)(0.857,0.125)(0.878,0.107)(0.898,0.089)(0.918,0.071)(0.939,0.054)(0.959,0.036)(0.980,0.018)(1.000,0.000)};
\draw[blue, opacity=0.4] plot coordinates {(0.000,1.000)(0.020,0.977)(0.041,0.955)(0.061,0.932)(0.082,0.910)(0.102,0.888)(0.122,0.865)(0.143,0.843)(0.163,0.821)(0.184,0.798)(0.204,0.776)(0.224,0.754)(0.245,0.732)(0.265,0.710)(0.286,0.689)(0.306,0.667)(0.327,0.645)(0.347,0.624)(0.367,0.603)(0.388,0.581)(0.408,0.560)(0.429,0.539)(0.449,0.518)(0.469,0.498)(0.490,0.477)(0.510,0.457)(0.531,0.437)(0.551,0.416)(0.571,0.396)(0.592,0.377)(0.612,0.357)(0.633,0.337)(0.653,0.318)(0.673,0.298)(0.694,0.279)(0.714,0.260)(0.735,0.241)(0.755,0.222)(0.776,0.203)(0.796,0.184)(0.816,0.166)(0.837,0.147)(0.857,0.129)(0.878,0.110)(0.898,0.092)(0.918,0.073)(0.939,0.055)(0.959,0.037)(0.980,0.018)(1.000,0.000)};
\draw[blue, opacity=0.4] plot coordinates {(0.000,1.000)(0.020,0.978)(0.041,0.956)(0.061,0.934)(0.082,0.912)(0.102,0.889)(0.122,0.867)(0.143,0.845)(0.163,0.823)(0.184,0.802)(0.204,0.780)(0.224,0.758)(0.245,0.736)(0.265,0.715)(0.286,0.693)(0.306,0.672)(0.327,0.650)(0.347,0.629)(0.367,0.608)(0.388,0.587)(0.408,0.566)(0.429,0.545)(0.449,0.524)(0.469,0.504)(0.490,0.483)(0.510,0.463)(0.531,0.442)(0.551,0.422)(0.571,0.402)(0.592,0.382)(0.612,0.362)(0.633,0.343)(0.653,0.323)(0.673,0.303)(0.694,0.284)(0.714,0.265)(0.735,0.245)(0.755,0.226)(0.776,0.207)(0.796,0.188)(0.816,0.169)(0.837,0.150)(0.857,0.131)(0.878,0.112)(0.898,0.094)(0.918,0.075)(0.939,0.056)(0.959,0.037)(0.980,0.019)(1.000,0.000)};
\draw[blue, opacity=0.4] plot coordinates {(0.000,1.000)(0.020,0.978)(0.041,0.956)(0.061,0.935)(0.082,0.913)(0.102,0.891)(0.122,0.869)(0.143,0.848)(0.163,0.826)(0.184,0.804)(0.204,0.783)(0.224,0.761)(0.245,0.740)(0.265,0.718)(0.286,0.697)(0.306,0.676)(0.327,0.654)(0.347,0.633)(0.367,0.612)(0.388,0.591)(0.408,0.571)(0.429,0.550)(0.449,0.529)(0.469,0.509)(0.490,0.488)(0.510,0.468)(0.531,0.447)(0.551,0.427)(0.571,0.407)(0.592,0.387)(0.612,0.367)(0.633,0.347)(0.653,0.327)(0.673,0.308)(0.694,0.288)(0.714,0.268)(0.735,0.249)(0.755,0.229)(0.776,0.210)(0.796,0.191)(0.816,0.172)(0.837,0.152)(0.857,0.133)(0.878,0.114)(0.898,0.095)(0.918,0.076)(0.939,0.057)(0.959,0.038)(0.980,0.019)(1.000,0.000)};
\draw[blue, opacity=0.4] plot coordinates {(0.000,1.000)(0.020,0.978)(0.041,0.957)(0.061,0.935)(0.082,0.914)(0.102,0.892)(0.122,0.871)(0.143,0.849)(0.163,0.828)(0.184,0.806)(0.204,0.785)(0.224,0.764)(0.245,0.742)(0.265,0.721)(0.286,0.700)(0.306,0.679)(0.327,0.658)(0.347,0.637)(0.367,0.616)(0.388,0.595)(0.408,0.574)(0.429,0.554)(0.449,0.533)(0.469,0.512)(0.490,0.492)(0.510,0.472)(0.531,0.451)(0.551,0.431)(0.571,0.411)(0.592,0.391)(0.612,0.371)(0.633,0.351)(0.653,0.331)(0.673,0.311)(0.694,0.291)(0.714,0.271)(0.735,0.252)(0.755,0.232)(0.776,0.213)(0.796,0.193)(0.816,0.174)(0.837,0.154)(0.857,0.135)(0.878,0.116)(0.898,0.096)(0.918,0.077)(0.939,0.058)(0.959,0.038)(0.980,0.019)(1.000,0.000)};
\draw[blue, opacity=0.4] plot coordinates {(0.000,1.000)(0.020,0.979)(0.041,0.957)(0.061,0.936)(0.082,0.915)(0.102,0.893)(0.122,0.872)(0.143,0.851)(0.163,0.829)(0.184,0.808)(0.204,0.787)(0.224,0.766)(0.245,0.745)(0.265,0.724)(0.286,0.703)(0.306,0.682)(0.327,0.661)(0.347,0.640)(0.367,0.619)(0.388,0.598)(0.408,0.577)(0.429,0.557)(0.449,0.536)(0.469,0.516)(0.490,0.495)(0.510,0.475)(0.531,0.455)(0.551,0.434)(0.571,0.414)(0.592,0.394)(0.612,0.374)(0.633,0.354)(0.653,0.334)(0.673,0.314)(0.694,0.294)(0.714,0.274)(0.735,0.254)(0.755,0.234)(0.776,0.215)(0.796,0.195)(0.816,0.176)(0.837,0.156)(0.857,0.136)(0.878,0.117)(0.898,0.097)(0.918,0.078)(0.939,0.058)(0.959,0.039)(0.980,0.019)(1.000,0.000)};
\draw[blue, opacity=0.4] plot coordinates {(0.000,1.000)(0.020,0.979)(0.041,0.958)(0.061,0.936)(0.082,0.915)(0.102,0.894)(0.122,0.873)(0.143,0.852)(0.163,0.831)(0.184,0.810)(0.204,0.789)(0.224,0.768)(0.245,0.747)(0.265,0.726)(0.286,0.705)(0.306,0.684)(0.327,0.663)(0.347,0.642)(0.367,0.621)(0.388,0.601)(0.408,0.580)(0.429,0.559)(0.449,0.539)(0.469,0.518)(0.490,0.498)(0.510,0.478)(0.531,0.457)(0.551,0.437)(0.571,0.417)(0.592,0.396)(0.612,0.376)(0.633,0.356)(0.653,0.336)(0.673,0.316)(0.694,0.296)(0.714,0.276)(0.735,0.256)(0.755,0.236)(0.776,0.217)(0.796,0.197)(0.816,0.177)(0.837,0.157)(0.857,0.138)(0.878,0.118)(0.898,0.098)(0.918,0.079)(0.939,0.059)(0.959,0.039)(0.980,0.020)(1.000,0.000)};
\draw[blue, opacity=0.4] plot coordinates {(0.000,1.000)(0.020,0.979)(0.041,0.958)(0.061,0.937)(0.082,0.916)(0.102,0.895)(0.122,0.874)(0.143,0.853)(0.163,0.832)(0.184,0.811)(0.204,0.790)(0.224,0.769)(0.245,0.748)(0.265,0.727)(0.286,0.706)(0.306,0.686)(0.327,0.665)(0.347,0.644)(0.367,0.623)(0.388,0.603)(0.408,0.582)(0.429,0.562)(0.449,0.541)(0.469,0.521)(0.490,0.500)(0.510,0.480)(0.531,0.459)(0.551,0.439)(0.571,0.419)(0.592,0.399)(0.612,0.378)(0.633,0.358)(0.653,0.338)(0.673,0.318)(0.694,0.298)(0.714,0.278)(0.735,0.258)(0.755,0.238)(0.776,0.218)(0.796,0.198)(0.816,0.178)(0.837,0.158)(0.857,0.138)(0.878,0.119)(0.898,0.099)(0.918,0.079)(0.939,0.059)(0.959,0.040)(0.980,0.020)(1.000,0.000)};
\draw[blue, opacity=0.4] plot coordinates {(0.000,1.000)(0.020,0.979)(0.041,0.958)(0.061,0.937)(0.082,0.916)(0.102,0.895)(0.122,0.874)(0.143,0.854)(0.163,0.833)(0.184,0.812)(0.204,0.791)(0.224,0.770)(0.245,0.749)(0.265,0.729)(0.286,0.708)(0.306,0.687)(0.327,0.666)(0.347,0.646)(0.367,0.625)(0.388,0.604)(0.408,0.584)(0.429,0.563)(0.449,0.543)(0.469,0.522)(0.490,0.502)(0.510,0.482)(0.531,0.461)(0.551,0.441)(0.571,0.421)(0.592,0.400)(0.612,0.380)(0.633,0.360)(0.653,0.340)(0.673,0.319)(0.694,0.299)(0.714,0.279)(0.735,0.259)(0.755,0.239)(0.776,0.219)(0.796,0.199)(0.816,0.179)(0.837,0.159)(0.857,0.139)(0.878,0.119)(0.898,0.099)(0.918,0.080)(0.939,0.060)(0.959,0.040)(0.980,0.020)(1.000,0.000)};
\draw[blue, opacity=0.4] plot coordinates {(0.000,1.000)(0.020,0.979)(0.041,0.958)(0.061,0.937)(0.082,0.917)(0.102,0.896)(0.122,0.875)(0.143,0.854)(0.163,0.833)(0.184,0.813)(0.204,0.792)(0.224,0.771)(0.245,0.750)(0.265,0.730)(0.286,0.709)(0.306,0.688)(0.327,0.668)(0.347,0.647)(0.367,0.626)(0.388,0.606)(0.408,0.585)(0.429,0.565)(0.449,0.544)(0.469,0.524)(0.490,0.503)(0.510,0.483)(0.531,0.463)(0.551,0.442)(0.571,0.422)(0.592,0.402)(0.612,0.381)(0.633,0.361)(0.653,0.341)(0.673,0.321)(0.694,0.301)(0.714,0.280)(0.735,0.260)(0.755,0.240)(0.776,0.220)(0.796,0.200)(0.816,0.180)(0.837,0.160)(0.857,0.140)(0.878,0.120)(0.898,0.100)(0.918,0.080)(0.939,0.060)(0.959,0.040)(0.980,0.020)(1.000,0.000)};
\draw[blue, opacity=0.4] plot coordinates {(0.000,1.000)(0.020,0.979)(0.041,0.958)(0.061,0.938)(0.082,0.917)(0.102,0.896)(0.122,0.875)(0.143,0.855)(0.163,0.834)(0.184,0.813)(0.204,0.793)(0.224,0.772)(0.245,0.751)(0.265,0.731)(0.286,0.710)(0.306,0.689)(0.327,0.669)(0.347,0.648)(0.367,0.628)(0.388,0.607)(0.408,0.586)(0.429,0.566)(0.449,0.546)(0.469,0.525)(0.490,0.505)(0.510,0.484)(0.531,0.464)(0.551,0.443)(0.571,0.423)(0.592,0.403)(0.612,0.383)(0.633,0.362)(0.653,0.342)(0.673,0.322)(0.694,0.302)(0.714,0.281)(0.735,0.261)(0.755,0.241)(0.776,0.221)(0.796,0.201)(0.816,0.181)(0.837,0.161)(0.857,0.140)(0.878,0.120)(0.898,0.100)(0.918,0.080)(0.939,0.060)(0.959,0.040)(0.980,0.020)(1.000,0.000)};
\draw[blue, opacity=0.4] plot coordinates {(0.000,1.000)(0.020,0.979)(0.041,0.959)(0.061,0.938)(0.082,0.917)(0.102,0.897)(0.122,0.876)(0.143,0.855)(0.163,0.834)(0.184,0.814)(0.204,0.793)(0.224,0.773)(0.245,0.752)(0.265,0.731)(0.286,0.711)(0.306,0.690)(0.327,0.670)(0.347,0.649)(0.367,0.628)(0.388,0.608)(0.408,0.587)(0.429,0.567)(0.449,0.547)(0.469,0.526)(0.490,0.506)(0.510,0.485)(0.531,0.465)(0.551,0.444)(0.571,0.424)(0.592,0.404)(0.612,0.383)(0.633,0.363)(0.653,0.343)(0.673,0.323)(0.694,0.302)(0.714,0.282)(0.735,0.262)(0.755,0.242)(0.776,0.222)(0.796,0.201)(0.816,0.181)(0.837,0.161)(0.857,0.141)(0.878,0.121)(0.898,0.101)(0.918,0.080)(0.939,0.060)(0.959,0.040)(0.980,0.020)(1.000,0.000)};
\draw[blue, opacity=0.4] plot coordinates {(0.000,1.000)(0.020,0.979)(0.041,0.959)(0.061,0.938)(0.082,0.917)(0.102,0.897)(0.122,0.876)(0.143,0.856)(0.163,0.835)(0.184,0.814)(0.204,0.794)(0.224,0.773)(0.245,0.752)(0.265,0.732)(0.286,0.711)(0.306,0.691)(0.327,0.670)(0.347,0.650)(0.367,0.629)(0.388,0.609)(0.408,0.588)(0.429,0.568)(0.449,0.547)(0.469,0.527)(0.490,0.506)(0.510,0.486)(0.531,0.466)(0.551,0.445)(0.571,0.425)(0.592,0.405)(0.612,0.384)(0.633,0.364)(0.653,0.344)(0.673,0.323)(0.694,0.303)(0.714,0.283)(0.735,0.263)(0.755,0.242)(0.776,0.222)(0.796,0.202)(0.816,0.182)(0.837,0.161)(0.857,0.141)(0.878,0.121)(0.898,0.101)(0.918,0.081)(0.939,0.061)(0.959,0.040)(0.980,0.020)(1.000,0.000)};
\draw[blue, opacity=0.4] plot coordinates {(0.000,1.000)(0.020,0.979)(0.041,0.959)(0.061,0.938)(0.082,0.918)(0.102,0.897)(0.122,0.876)(0.143,0.856)(0.163,0.835)(0.184,0.815)(0.204,0.794)(0.224,0.774)(0.245,0.753)(0.265,0.732)(0.286,0.712)(0.306,0.691)(0.327,0.671)(0.347,0.650)(0.367,0.630)(0.388,0.609)(0.408,0.589)(0.429,0.568)(0.449,0.548)(0.469,0.528)(0.490,0.507)(0.510,0.487)(0.531,0.466)(0.551,0.446)(0.571,0.426)(0.592,0.405)(0.612,0.385)(0.633,0.365)(0.653,0.344)(0.673,0.324)(0.694,0.304)(0.714,0.283)(0.735,0.263)(0.755,0.243)(0.776,0.222)(0.796,0.202)(0.816,0.182)(0.837,0.162)(0.857,0.142)(0.878,0.121)(0.898,0.101)(0.918,0.081)(0.939,0.061)(0.959,0.040)(0.980,0.020)(1.000,0.000)};
\draw[blue, opacity=0.4] plot coordinates {(0.000,1.000)(0.020,0.979)(0.041,0.959)(0.061,0.938)(0.082,0.918)(0.102,0.897)(0.122,0.877)(0.143,0.856)(0.163,0.835)(0.184,0.815)(0.204,0.794)(0.224,0.774)(0.245,0.753)(0.265,0.733)(0.286,0.712)(0.306,0.692)(0.327,0.671)(0.347,0.651)(0.367,0.630)(0.388,0.610)(0.408,0.589)(0.429,0.569)(0.449,0.549)(0.469,0.528)(0.490,0.508)(0.510,0.487)(0.531,0.467)(0.551,0.446)(0.571,0.426)(0.592,0.406)(0.612,0.385)(0.633,0.365)(0.653,0.345)(0.673,0.324)(0.694,0.304)(0.714,0.284)(0.735,0.263)(0.755,0.243)(0.776,0.223)(0.796,0.203)(0.816,0.182)(0.837,0.162)(0.857,0.142)(0.878,0.121)(0.898,0.101)(0.918,0.081)(0.939,0.061)(0.959,0.040)(0.980,0.020)(1.000,0.000)};
\draw[blue, opacity=0.4] plot coordinates {(0.000,1.000)(0.020,0.979)(0.041,0.959)(0.061,0.938)(0.082,0.918)(0.102,0.897)(0.122,0.877)(0.143,0.856)(0.163,0.836)(0.184,0.815)(0.204,0.795)(0.224,0.774)(0.245,0.754)(0.265,0.733)(0.286,0.713)(0.306,0.692)(0.327,0.672)(0.347,0.651)(0.367,0.631)(0.388,0.610)(0.408,0.590)(0.429,0.569)(0.449,0.549)(0.469,0.529)(0.490,0.508)(0.510,0.488)(0.531,0.467)(0.551,0.447)(0.571,0.427)(0.592,0.406)(0.612,0.386)(0.633,0.365)(0.653,0.345)(0.673,0.325)(0.694,0.304)(0.714,0.284)(0.735,0.264)(0.755,0.243)(0.776,0.223)(0.796,0.203)(0.816,0.183)(0.837,0.162)(0.857,0.142)(0.878,0.122)(0.898,0.101)(0.918,0.081)(0.939,0.061)(0.959,0.041)(0.980,0.020)(1.000,0.000)};
\draw[blue, opacity=0.4] plot coordinates {(0.000,1.000)(0.020,0.979)(0.041,0.959)(0.061,0.938)(0.082,0.918)(0.102,0.897)(0.122,0.877)(0.143,0.856)(0.163,0.836)(0.184,0.815)(0.204,0.795)(0.224,0.774)(0.245,0.754)(0.265,0.733)(0.286,0.713)(0.306,0.692)(0.327,0.672)(0.347,0.652)(0.367,0.631)(0.388,0.611)(0.408,0.590)(0.429,0.570)(0.449,0.549)(0.469,0.529)(0.490,0.508)(0.510,0.488)(0.531,0.468)(0.551,0.447)(0.571,0.427)(0.592,0.407)(0.612,0.386)(0.633,0.366)(0.653,0.345)(0.673,0.325)(0.694,0.305)(0.714,0.284)(0.735,0.264)(0.755,0.244)(0.776,0.223)(0.796,0.203)(0.816,0.183)(0.837,0.162)(0.857,0.142)(0.878,0.122)(0.898,0.102)(0.918,0.081)(0.939,0.061)(0.959,0.041)(0.980,0.020)(1.000,0.000)};
\draw[blue, opacity=0.4] plot coordinates {(0.000,1.000)(0.020,0.980)(0.041,0.959)(0.061,0.939)(0.082,0.918)(0.102,0.898)(0.122,0.877)(0.143,0.857)(0.163,0.836)(0.184,0.816)(0.204,0.795)(0.224,0.775)(0.245,0.754)(0.265,0.734)(0.286,0.713)(0.306,0.693)(0.327,0.672)(0.347,0.652)(0.367,0.631)(0.388,0.611)(0.408,0.590)(0.429,0.570)(0.449,0.550)(0.469,0.529)(0.490,0.509)(0.510,0.488)(0.531,0.468)(0.551,0.448)(0.571,0.427)(0.592,0.407)(0.612,0.386)(0.633,0.366)(0.653,0.346)(0.673,0.325)(0.694,0.305)(0.714,0.285)(0.735,0.264)(0.755,0.244)(0.776,0.224)(0.796,0.203)(0.816,0.183)(0.837,0.163)(0.857,0.142)(0.878,0.122)(0.898,0.102)(0.918,0.081)(0.939,0.061)(0.959,0.041)(0.980,0.020)(1.000,0.000)};
\draw[blue, opacity=0.4] plot coordinates {(0.000,1.000)(0.020,0.980)(0.041,0.959)(0.061,0.939)(0.082,0.918)(0.102,0.898)(0.122,0.877)(0.143,0.857)(0.163,0.836)(0.184,0.816)(0.204,0.795)(0.224,0.775)(0.245,0.754)(0.265,0.734)(0.286,0.713)(0.306,0.693)(0.327,0.672)(0.347,0.652)(0.367,0.632)(0.388,0.611)(0.408,0.591)(0.429,0.570)(0.449,0.550)(0.469,0.529)(0.490,0.509)(0.510,0.489)(0.531,0.468)(0.551,0.448)(0.571,0.427)(0.592,0.407)(0.612,0.387)(0.633,0.366)(0.653,0.346)(0.673,0.326)(0.694,0.305)(0.714,0.285)(0.735,0.264)(0.755,0.244)(0.776,0.224)(0.796,0.203)(0.816,0.183)(0.837,0.163)(0.857,0.142)(0.878,0.122)(0.898,0.102)(0.918,0.081)(0.939,0.061)(0.959,0.041)(0.980,0.020)(1.000,0.000)};
\draw[blue, opacity=0.4] plot coordinates {(0.000,1.000)(0.020,0.980)(0.041,0.959)(0.061,0.939)(0.082,0.918)(0.102,0.898)(0.122,0.877)(0.143,0.857)(0.163,0.836)(0.184,0.816)(0.204,0.795)(0.224,0.775)(0.245,0.754)(0.265,0.734)(0.286,0.714)(0.306,0.693)(0.327,0.673)(0.347,0.652)(0.367,0.632)(0.388,0.611)(0.408,0.591)(0.429,0.571)(0.449,0.550)(0.469,0.530)(0.490,0.509)(0.510,0.489)(0.531,0.468)(0.551,0.448)(0.571,0.428)(0.592,0.407)(0.612,0.387)(0.633,0.366)(0.653,0.346)(0.673,0.326)(0.694,0.305)(0.714,0.285)(0.735,0.265)(0.755,0.244)(0.776,0.224)(0.796,0.204)(0.816,0.183)(0.837,0.163)(0.857,0.142)(0.878,0.122)(0.898,0.102)(0.918,0.081)(0.939,0.061)(0.959,0.041)(0.980,0.020)(1.000,0.000)};
\draw[blue, opacity=0.4] plot coordinates {(0.000,1.000)(0.020,0.980)(0.041,0.959)(0.061,0.939)(0.082,0.918)(0.102,0.898)(0.122,0.877)(0.143,0.857)(0.163,0.836)(0.184,0.816)(0.204,0.795)(0.224,0.775)(0.245,0.755)(0.265,0.734)(0.286,0.714)(0.306,0.693)(0.327,0.673)(0.347,0.652)(0.367,0.632)(0.388,0.612)(0.408,0.591)(0.429,0.571)(0.449,0.550)(0.469,0.530)(0.490,0.509)(0.510,0.489)(0.531,0.469)(0.551,0.448)(0.571,0.428)(0.592,0.407)(0.612,0.387)(0.633,0.367)(0.653,0.346)(0.673,0.326)(0.694,0.305)(0.714,0.285)(0.735,0.265)(0.755,0.244)(0.776,0.224)(0.796,0.204)(0.816,0.183)(0.837,0.163)(0.857,0.143)(0.878,0.122)(0.898,0.102)(0.918,0.081)(0.939,0.061)(0.959,0.041)(0.980,0.020)(1.000,0.000)};
\draw[blue, opacity=0.4] plot coordinates {(0.000,1.000)(0.020,0.980)(0.041,0.959)(0.061,0.939)(0.082,0.918)(0.102,0.898)(0.122,0.877)(0.143,0.857)(0.163,0.836)(0.184,0.816)(0.204,0.796)(0.224,0.775)(0.245,0.755)(0.265,0.734)(0.286,0.714)(0.306,0.693)(0.327,0.673)(0.347,0.652)(0.367,0.632)(0.388,0.612)(0.408,0.591)(0.429,0.571)(0.449,0.550)(0.469,0.530)(0.490,0.510)(0.510,0.489)(0.531,0.469)(0.551,0.448)(0.571,0.428)(0.592,0.408)(0.612,0.387)(0.633,0.367)(0.653,0.346)(0.673,0.326)(0.694,0.306)(0.714,0.285)(0.735,0.265)(0.755,0.244)(0.776,0.224)(0.796,0.204)(0.816,0.183)(0.837,0.163)(0.857,0.143)(0.878,0.122)(0.898,0.102)(0.918,0.081)(0.939,0.061)(0.959,0.041)(0.980,0.020)(1.000,0.000)};
\draw[very thick, red] (0,1) -- (1,0);
\end{tikzpicture}

    \caption{Изменение давления вдоль пласта с течением времени}
    \label{fig:pressure_evolution}
\end{figure}

\newpage
\section*{ЗАКЛЮЧЕНИЕ}
\addcontentsline{toc}{section}{ЗАКЛЮЧЕНИЕ}

\newpage
\section*{ПРИЛОЖЕНИЕ}
\addcontentsline{toc}{section}{ПРИЛОЖЕНИЕ}

\newpage
\section*{СПИСОК ЛИТЕРАТУРЫ}
\addcontentsline{toc}{section}{СПИСОК ЛИТЕРАТУРЫ}

\begingroup
\renewcommand{\refname}{}

\begin{thebibliography}{99}

% --- Литература по гидромеханике и фильтрации ---
\bibitem{kan} Каневская Р. Д. Математическое моделирование гидродинамических процессов разработки месторождений углеводородов. --- Москва--Ижевск: Институт компьютерных исследований, 2002. --- 140 с.

\bibitem{aziz} Азиз Х., Сеттари Э. Математическое моделирование пластовых систем. --- Москва--Ижевск: Институт компьютерных исследований, 2004. --- 416 с.

\bibitem{afan1} Афанаскин И. В., Кац Р. М., Королев А. В. Моделирование разработки нефтяных месторождений: курс лекций. --- Сургут--Москва: СурГУ, ФГУ ФНЦ НИИСИ РАН, 2015. --- 87 с.

\bibitem{afan2} Афанаскин И. В., Юдин В. А., Вольпин С. Г. Подземная гидромеханика: курс лекций. --- Сургут--Москва: СурГУ, ФГУ ФНЦ НИИСИ РАН, 2015. --- 100 с.

\bibitem{kvesko} Квеско Б. Б. Подземная гидромеханика: учебное пособие. --- Томск: Изд-во ТПУ, 2010. --- 181 с.

% --- Зарубежные статьи и обзоры ---
\bibitem{khalili} Khalili Y., Ahmadi M. Reservoir Modeling \& Simulation: Advancements, Challenges, and Future Perspectives // Journal of Chemical and Petroleum Engineering. --- 2023. --- Vol. 57, No. 2. --- P. 343--364.

% --- Литература по PINN ---
\bibitem{bombini} Bombini A. Lecture notes on Physics Informed Neural Networks, Neural Operators, and their applications. --- Free University of Bozen-Bolzano, 2026.

\bibitem{baty} Baty H. A hands-on introduction to Physics-Informed Neural Networks for solving partial differential equations with benchmark tests taken from astrophysics and plasma physics. --- 2024. --- hal-04491808.

\bibitem{navarro} Navarro L. M., Martin-Moreno L., Rodrigo S. G. Solving differential equations with Deep Learning: a beginner's guide // arXiv:2307.11237. --- 2023.

% --- Практические руководства и код ---
% \bibitem{shaf} SHAFNehal. Learning-PINNs-in-Pytorch-Physics-Informed-Machine-Learning: GitHub repository. --- 2025.

% \bibitem{dedbin} dedbin. PINN Heat1D: GitHub repository. --- 2025.

\end{thebibliography}

\end{document}